\section*{THE THEORY OF SETS.}
\phantomsection
\addcontentsline{toc}{section}{THE THEORY OF SETS.}

It can be said that at all times, mathematicians and philosophers have used reasoned arguments from the theory of sets in a more or less conscious way; but in the history of their conceptions of this subject, one must separate sharply all questions linked to the idea of cardinal number (and in particular to the notion of infinity) from those that only introduce the notions of belonging and inclusion. These latter are among the most intuitive and seem never to have raised controversy: it is on them that one can most easily base a theory of syllogism (as Leibniz and Euler were to show us), or axioms such as ``the whole is greater than the sum of its parts'', without talking about that which, in geometry, is concerned with intersections of curves and surfaces. Until the end of the XIXth century, no problem arises in talking about the set (or ``class'' with some authors) of objects possessing such and such a given property; \(^{38}\) and the famous ``definition'' given by Cantor (``By a set is meant a gathering into one whole of objects which are quite distinct in our intuition or our thought'' ({[}47{]}, p.~282)) will give rise, at the time of its publication, to hardly any objections. \(^{39}\) It is altogether different as soon as the notion of set is mixed with that of number or quantity. The question of the infinite divisibility of the line (doubtless posed already by the first Pythagoreans) was, as is known, to lead to considerable philosophical difficulties: from the Eleates to Bolzano and Cantor, mathematicians and philosophers will throw themselves without success against the paradox of the finite quantity made up of an infinity of points without size. It would be of no interest to us to retrace, even summarily, the interminable and impassioned polemics that are aroused by this problem, that constituted a ground particularly favourable to metaphysical or theological digressions; let us note only the point of view at which, from Antiquity, most mathematicians stop. It consists essentially in refusing to debate, short of being able to conclude in an irrefutable way --- an attitude that we will find again among the modern formalists; in the same way as these latter arrange for the elimination of the intervention of ``paradoxical'' sets (see below pp.~31 ff.), the classical mathematicians carefully avoid introducing into their reasoning the ``actual infinity'' (that is to say sets made up of an infinity of objects conceived as existing simultaneously, at least in their thought), and content themselves with ``potential infinity'', that is to say the possibility of increasing every given quantity (or of decreasing it if it is a case of a ``continuous'' quantity). \(^{40}\) If this point of view contained a certain dose of hypocrisy, \(^{41}\) it allowed however the development of the larger part of classical mathematics (including the theory of ratios and later the infinitesimal Calculus); \(^{42}\) it even seemed an excellent safeguard, above all after the quarrels caused by the infinitely small, and had become a dogma almost universally subscribed to until well into the XIXth century.

A first glimmer of the general notion of equipotence appears in a remark of Galileo ({[}122 b{]}, v. VIII, pp.~78-80): he observes that the map \(n \mapsto n^2\) establishes a bijective correspondence between the natural numbers and their squares and that it follows that the axiom ``the whole is greater than the part''

could not be applied to an infinite set. But far from inaugurating a rational study of infinite sets, this remark seems to have had no other effect than to reinforce mistrust of actual infinity; it is already the conclusion of Galileo himself, and Cauchy, in 1833, only quotes him in order to approve of his attitude.

The necessities of Analysis --- and particularly the deep study of functions of real variables, which was taking place during the whole of the XIXth century --- are at the origins of what was to become the modern theory of sets. When Bolzano, in 1817, proves the existence of the lower limit of a set bounded below in R {[}27 c{]}, he still reasons, like the greater part of his contemporaries, ``by comprehension'', speaking not of an arbitrary set of real numbers, but of an arbitrary property of these latter. But when, 30 years later, he writes his ``Paradoxien des Unendlichen'' {[}27 b{]} (published in 1851, three years after his death), he does not hesitate to claim the right to the existence of the ``actual infinity'' and to speak of arbitrary sets. He defines, in this work, the general notion of the equipotence of two sets, and proves that two compact intervals in R are equipotent; he notes also that the characteristic difference between finite sets and infinite sets consists in that an infinite set E is equipotent to a subset distinct from E, but he does not give any convincing proof of this assertion. The general tone of this work is in any case more philosophical than mathematical, and for lack of disassociating sufficiently clearly the notion of the power of a set from that of size or order of infinity, Bolzano fails in his attempts to construct infinite sets of greater and greater power, and lets himself be drawn on this occasion into mixing his reasoning with considerations on divergent series that are totally devoid of sense.

It is due to the genius of G. Cantor that the theory of sets as we know it today was created. He also started with Analysis, and his work on trigonometric series, inspired by those of Riemann, bring him naturally, in 1872, to a first attempt to classify ``exceptional'' sets which occur in this theory, \(^{43}\) by means of the notion of successive ``derived sets'', which he introduces at this point. It is no doubt in connection with these researches, and also with his method for defining the real numbers ({[}47{]}, pp.~92-97), that Cantor begins to be interested in the problems of equipotence, for in 1873, he remarks that the set of rational numbers (or the set of algebraic numbers) is countable; and in his correspondence with Dedekind, which starts about this time {[}48{]}, he is seen posing the question of the equipotence between the set of whole numbers and the set of real numbers, which he succeeds in solving negatively a few weeks later. Then, from 1874 on, it is the problem of dimension that preoccupies him, and for three years he seeks in vain to establish the impossibility of a bijective correspondence between R and \(R^{n}(n > 1)\) , before succeeding, to his great surprise, \(^{44}\) in defining such a correspondence. Having such results as new as they were surprising, he devotes himself entirely to the theory of sets. In a series of 6 memoirs published in the Mathematische Annalen between 1878 and 1884, he tackles both the problems of equipotence, the theory of totally ordered sets, the topological properties of R and of R \(^{n}\) and the problem of measure; and it is admirable to see with what clarity these notions appear in his hands little by little, notions which appear so inextricably tied to the classical conception of the ``continuum''. From 1880, he has the idea of iterating ``transfinitely'' the construction of ``derived sets''; but this idea takes on life only two years later, with the introduction of well ordered sets, one of the most original discoveries of Cantor, thanks to which he can tackle a detailed study of cardinal numbers and formulate the ``problem of the continuum'' {[}47{]}.

It was impossible that such bold conceptions, reversing a twice millennial tradition, and leading to results so unexpected and of such a paradoxical appearance, would be accepted without resistance. In fact, amongst the mathematicians who were influential at that time in Germany, Weierstrass was the only one to follow with some favour the work of Cantor (his erstwhile pupil); this latter was to run up on the other hand against the irreducible opposition of Schwarz and especially of Kronecker. \(^{45}\) It is, it would appear, as much the constant tension generated by the opposition to his ideas, as his fruitless efforts to prove the continuum hypothesis, which caused in Cantor the first symptoms of a nervous illness from which his mathematical productivity would suffer. \(^{46}\) He only really took up again his interest in the theory of sets in 1887, and his last publications date from 1895-1897; he develops there mostly the theory of totally ordered sets and the calculus of ordinals. He had also proved in 1890 the inequality \(m < 2^{m}\) ; however, not only did the problem of the continuum remain without answer, but there remained in the theory of cardinals a more serious gap, for Cantor had been unable to establish the existence of a well ordering between arbitrary cardinals. This gap was going to be filled, on the one hand by the theorem of F. Bernstein (1897) showing that the relations \(a \leq b\) and \(b \leq a\) imply \(a = b\) , \(^{47}\) and above all by the theorem of Zermelo {[}342 a{]} proving the existence of a well order on every set --- a theorem conjectured already in 1883 by Cantor ({[}47{]}, p.~169).

However Dedekind, from the beginning, had not ceased to follow with a sustained interest the researches of Cantor; but while this latter concentrated his attention on infinite sets and their classification, Dedekind followed his own thoughts on the notion of number (which had already led him to his definition of irrationals by ``cuts''). In his opuscule Was sind und was sollen die Zahlen, published in 1888, but the essentials dating from 1872-78 ({[}79{]}, v. III, p.~335), he shows how the notion of natural number (on which, as has been seen, the whole of classical Mathematics had ended up by relying) could itself be derived from fundamental notions in the theory of sets. Developing (no doubt the first explicitly) the elementary properties of arbitrary maps from one set to another (neglected until Cantor, who was only interested in bijective correspondences), he introduces, for every map f of a set E into itself, the notion of ``chain'' of an element \(a \in E\) relative to f, namely the intersection of the sets \(K \subset E\) such that \(a \in K\) and \(f(K) \subset K\) .⁴⁸ He then takes as his definition of an infinite set E the fact that there exists a bijective map \(\phi\) of E into E such that \(\phi(E) \neq E\) ;⁴⁹ if further there exists such a map \(\phi\) and an element \(a \notin \phi(E)\) for which E is the chain of a, Dedekind says that E is ``simply infinite'', remarks that the ``axioms of Peano'' are then satisfied and shows (before Peano) how starting from there all the theorems of elementary arithmetic are obtained. The only missing part of his exposé is the axiom of infinity that Dedekind (following Bolzano) believes himself able to prove by considering the ``world of thought'' of humans (``Gedankenwelt'') as a set.⁵⁰

On another tack, Dedekind had been led by his work on arithmetic (and notably the theory of ideals) to consider the notion of ordered set in a more general setting than Cantor. While this latter restricts himself exclusively to totally ordered sets, \(^{51}\) Dedekind tackles the general case, and makes in particular a deep study of lattices ({[}79{]}, v. II, pp.~236-271). This work was hardly noticed at the time; although these results, found again by several authors, have been the object of numerous publications since 1935, their historical importance arises much less from the possibilities for applications, fairly thin no doubt, of this theory, than from the fact that they constituted one of the first examples of careful axiomatic construction. On the other hand, the first results of Cantor on sets which were countable or had the power of the continuum were rapidly to have numerous and important applications right up to the most classical questions of Analysis \(^{52}\) (without talking naturally of those parts of the Cantorian work that inaugurated General Topology and measure theory; on that see pp.~141 and 221). Moreover, from the last years of the XIXth century the first uses of the principle of transfinite induction were appearing, having become, especially after the proof of the theorem of Zermelo, an indispensable tool in all parts of modern mathematics. Kuratowski was to give, in 1922, an often more amenable version of this principle, avoiding the use of well ordered sets ({[}189 a{]}, p.~89); it is in this form, found again later by Zorn {[}344{]}, that it is mainly used at the present time. \(^{53}\)

Towards the end of the XIXth century, the essential concepts of Cantor had thus gained their cause. \(^{54}\) We have seen that, at about this time, the formalisation of mathematics is achieved and the use of the axiomatic method is almost universally agreed. But, at the same time, a ``crisis of foundations'' of a rare violence was opening up, which was to shake the mathematical world for more than 30 years, and to seem at times to compromise, not only all the recent acquisitions, but even the most classical parts of Mathematics.

